\chapter{补充}

\section{数值格式整理}

\subsection*{一维常系数线性对流方程}

考虑方程 $u_t + a u_x = 0$，其中 $a \neq 0$ 为常数，常见的数值格式如下：
\begin{itemize}
    \item 迎风格式（FTFS / FTBS）
          \[
              \frac{v_j^{n+1} - v_j^n}{\Delta t} + a \frac{v_{j+1}^n - v_{j-1}^{n}}{2\Delta x} = \frac{|a|}{2\Delta x} (v_{j+1}^n - 2 v_j^n + v_{j-1}^{n}).
          \]
    \item FTCS（反面示例，无条件不稳定）
          \[
              \frac{v_{j}^{n+1} - v_{j}^n}{\Delta t} + a \frac{v_{j+1}^n - v_{j-1}^n}{2\Delta x} = 0.
          \]
    \item Lax-Friedrichs 格式
          \begin{gather*}
              \frac{v_j^{n+1}-v_j^n}{\Delta t} + a \frac{v_{j+1}^n-v_{j-1}^n}{2\Delta x} ={} \left(\frac{\Delta x^2}{2\Delta t}\right) \frac{v_{j+1}^n-2 v_j^n+v_{j-1}^n}{\Delta x^2}     \\
              \frac{v_j^{n+1}-\frac12(v_{j-1}^n+v_{j+1}^n)}{\Delta t} + a \frac{v_{j+1}^n-v_{j-1}^n}{2\Delta x} ={} 0.
          \end{gather*}
    \item Lax-Wendroff 格式
          \begin{gather*}
              \frac{v_j^{n+1}-v_j^n}{\Delta t} + a \frac{v_{j+1}^n-v_{j-1}^n}{2\Delta x} = \left(\frac{a^2 \Delta t}{2}\right) \frac{v_{j+1}^n-2 v_j^n+v_{j-1}^n}{\Delta x^2} \\
              v_{j}^{n+1} = v_j^n + \Delta t \,(-a)\,\left(\frac{v_{j+1}^n- v_{j-1}^n}{2\Delta x}\right)
              + \frac{\Delta t^2}2\,a^2\,\left(\frac{v_{j+1}^n-2v_{j}^n+v_{j-1}^n}{\Delta x^2}\right).
          \end{gather*}
    \item 蛙跳格式（CTCS）
          \[
              \frac{v_{j}^{n+1} - v_{j}^{n-1}}{2\Delta t} + a \frac{v_{j+1}^n - v_{j-1}^n}{2\Delta x} = 0.
          \]
    \item 盒子格式（Wendroff）
          \begin{gather*}
              \frac{\frac12(v_j^{n+1} + v_{j+1}^{n+1}) - \frac12(v_j^n + v_{j+1}^n)}{\Delta t} + a \frac{\frac12(v_{j+1}^n + v_{j+1}^{n+1})-\frac12(v_j^n + v_j^{n+1})}{\Delta x} = 0 \\
              \left(1 - \frac{a \Delta t}{\Delta x}\right)v_j^{n+1} + \left(1 + \frac{a \Delta t}{\Delta x}\right)v_{j+1}^{n+1} =
              \left(1 + \frac{a \Delta t}{\Delta x}\right)v_j^{n} + \left(1 - \frac{a \Delta t}{\Delta x}\right)v_{j+1}^{n}.
          \end{gather*}
\end{itemize}


\subsection*{一维常系数线性扩散方程}

考虑方程 $u_t = b u_{xx}$，其中 $b > 0$ 为常数，常见的数值格式如下：
\begin{itemize}
    \item FTCS
          \[
              \frac{v_{j}^{n+1} - v_{j}^n}{\Delta t} = b \frac{v_{j+1}^n - 2v_{j}^n + v_{j-1}^n}{\Delta x^2}.
          \]
    \item BTCS
          \[
              \frac{v_{j}^{n+1} - v_{j}^n}{\Delta t} = b \frac{v_{j+1}^{n+1} - 2v_{j}^{n+1} + v_{j-1}^{n+1}}{\Delta x^2}.
          \]
    \item Richardson 格式（CTCS，反面示例，无条件不稳定）
          \[
              \frac{v_{j}^{n+1} - v_{j}^{n-1}}{2\Delta t} = b \frac{v_{j+1}^{n} - 2v_{j}^{n} + v_{j-1}^{n}}{\Delta x^2}.
          \]
    \item Du Fort-Frankel 格式
          \[
              \frac{v_j^{n+1}-v_{j-1}^n}{2\Delta t} = b \frac{v_{j+1}^n - (v_{j}^{n-1}+v_j^{n+1}) + v_{j-1}^n}{\Delta x^2}.
          \]
    \item $\theta$ 格式
          \[
              \frac{v_{j}^{n+1} - v_{j}^n}{\Delta t}
              =
              b \left[
                  \theta \frac{v_{j+1}^{n+1} - 2v_{j}^{n+1} + v_{j-1}^{n+1}}{\Delta x^2}
                  + (1-\theta) \frac{v_{j+1}^n - 2v_{j}^n + v_{j-1}^n}{\Delta x^2}
                  \right].
          \]
    \item Crank-Nicolson 格式：$\theta$ 格式取 $\theta = \frac12$
    \item Douglas 格式：$\theta$ 格式取 $\theta$ 满足
          \[
              \frac{b \Delta t}{\Delta x^2} = \frac{1}{6(1-2\theta)}.
          \]
    \item BDF 格式
          \[
              (1+\theta)\frac{v_{j}^{n+1} - v_{j}^{n}}{\Delta t}
              - \theta \frac{v_j^n - v_j^{n-1}}{\Delta t}
              = b \frac{v_{j+1}^{n} - 2v_{j}^{n} + v_{j-1}^{n}}{\Delta x^2}.
          \]
    \item Richtmyer 格式：BDF 格式取 $\theta = \frac12$
\end{itemize}

\subsection*{一维常系数线性对流方程的稳定性结论}

在纯初值或周期边界条件下，使用Fourier方法分析 $u_t+au_x=0$ 的几类经典格式的稳定性，结论如下（记 $r = \frac{a\Delta t}{\Delta x} \neq 0$）:
按照格式分类：
\begin{table}[H]
    \centering
    \begin{tabular}{c|c|c|c}
        \hline
        FTFS & $r(1+r) \le 0$           & BTFS & $r(1-r) \le 0$          \\
        \hline
        FTBS & $r(1-r) \ge 0$           & BTBS & $r(1+r) \ge 0$          \\
        \hline
        FTCS & 无条件不稳定（$r^2 \le 0$ 无法成立） & BTCS & 无条件稳定（$r^2 \ge 0$ 始终成立） \\
        \hline
    \end{tabular}
\end{table}
对于单侧差商格式，按照 $a$ 的正负分类：
\begin{enumerate}
    \item $a>0$：风向从左往右
          \begin{itemize}
              \item FTFS 不稳定
              \item FTBS 有条件稳定（$0< r \le 1$）
              \item BTFS 有条件稳定（$r \ge 1$）
              \item BTBS 无条件稳定
          \end{itemize}
    \item $a<0$：风向从右往左
          \begin{itemize}
              \item FTFS 有条件稳定（$-1 \le r < 0$）
              \item FTBS 不稳定
              \item BTFS 无条件稳定
              \item BTBS 有条件稳定（$r \le -1$）
          \end{itemize}
\end{enumerate}

\section{数值微分公式的 Mathematica 代码}

对于指定模板的数值微分公式（$a,m,n \ge 0$）
\[
    f^{(a)}(x) \approx \sum_{i=-m}^n c_i f(x + i h).
\]
具体系数 $\{c_i\}$ 需要通过联列线性方程组进行求解，计算过程非常无聊，可以通过 Mathematica 简化计算，具体代码如下（Mathematica 14.0）
\begin{lstlisting}
FDFormula[m_, n_, a_] :=
 Module[{mat, b, coeff, expr, i, j, jmax, f, x, h}, f = Unique[f];
  x = Unique[x]; h = Unique[h];
  mat = Table[If[j == 0, 1, i^j], {j, 0, m + n}, {i, -m, n}];
  b = Table[If[i == a, a!, 0], {i, 0, n + m}];
  coeff = LinearSolve[mat, b];
  expr = Sum[coeff[[m + 1 + i]] f[x + i h]/h^a, {i, -m, n}];
  For[j = a + 1, j <= m + n + 2, j++,
   If[Simplify@Sum[coeff[[m + 1 + i]]*i^j, {i, -m, n}] != 0, Break[]]];
  jmax = j;
  Module[{fout, xout, hout}, fout = Symbol["f"];
   xout = Symbol["x"]; hout = Symbol["h"];
   Print["coeff: ", coeff, "\nexpression: ",
    Format[expr /. {f -> fout, x -> xout, h -> hout},
     TraditionalForm], "\nseries: ",
    Series[expr /. {f -> fout, x -> xout, h -> hout}, {hout, 0,
      jmax - a}], "\norder: ", jmax - a];
   Print[Format[expr /. {f -> fout, x -> xout, h -> hout}, TeXForm],
    " = ", Format[
     Series[expr /. {f -> fout, x -> xout, h -> hout}, {hout, 0,
       jmax - a}], TeXForm]];];
  Return[coeff]]
\end{lstlisting}

\begin{example}
    使用 $x-h,x,x+h$ 三个点的中心模板求一阶导：
    \begin{lstlisting}[numbers=none]
FDFormula[1, 1, 1]

coeff: {-(1/2),0,1/2}
expression: f(h+x)/(2 h)-f(x-h)/(2 h)
series: (f^\[Prime])[x]+1/6 (f^(3))[x] h^2+O[h]^3
order: 2
    \end{lstlisting}
    对应数值微分公式
    \[
        \frac{f(h+x)}{2 h}-\frac{f(x-h)}{2 h}=f'(x)+\frac{1}{6} h^2 f^{(3)}(x)+O\left(h^3\right).
    \]
\end{example}


\begin{example}
    使用 $x,x+h,x+2h,x+3h$ 四个点的偏心模板求一阶导：
    \begin{lstlisting}[numbers=none]
FDFormula[0, 3, 1]

coeff: {-(11/6),3,-(3/2),1/3}
expression: -((11 f(x))/(6 h))+(3 f(h+x))/h-(3 f(2 h+x))/(2 h)+f(3 h+x)/(3 h)
series: (f^\[Prime])[x]+1/4 (f^(4))[x] h^3+O[h]^4
order: 3
    \end{lstlisting}
    对应数值微分公式
    \[
        -\frac{11 f(x)}{6 h}+\frac{3 f(h+x)}{h}-\frac{3 f(2 h+x)}{2 h}+\frac{f(3 h+x)}{3 h}=f'(x)+\frac{1}{4} h^3 f^{(4)}(x)+O\left(h^4\right).
    \]
\end{example}


\section{MPDE 方法的 Mathematica 代码}

MPDE 方法是一类用来分析数值格式的耗散色散性质的方法，它的思路为：
对于 $u_t = \mathcal{L}(u)$ ，我们通过离散设计得到了数值格式，
现在反过来，找出一个与数值格式最契合的 $u_t = \widetilde{\mathcal{L}}(u)$，通常是在 $u_t = \mathcal{L}(u)$ 的基础上加上一些依赖网格尺度的高阶余项，直接分析 $u_t = \widetilde{\mathcal{L}}(u)$ 的耗散和色散性质，尤其关注它多出来的高阶项，就可以得到数值格式的耗散和色散性质。

从差分格式找到对应的 $u_t = \widetilde{\mathcal{L}}(u)$ 是 MPDE 方法的核心，也是最繁琐的部分，基本做法就是将精确解在差分格式上进行泰勒展开，对其不断求导并与自身运算，消去除了 $u_t$ 之外的低价时间导，将时间导逐步转换为空间导，在分析时舍弃高阶项。
可以利用 Mathematica 简化符号计算过程，实现思路如下：
\begin{enumerate}
    \item 按照如下位置排序将对应项的系数存储为下三角矩阵
          \[
              \begin{bmatrix}
                  u                                              \\
                  u_t     & u_x                                  \\
                  u_{tt}  & u_{tx}  & u_{xx}                     \\
                  u_{ttt} & u_{ttx} & u_{txx} & u_{xxx}          \\
                  \vdots  & \vdots  & \vdots  & \vdots  & \ddots \\
              \end{bmatrix}.
          \]
          实际只使用固定尺寸 $n \times n$ 的矩阵来存储，对于更高阶系数直接忽略。
          例如初始状态的系数矩阵可以通过对差分格式的时间项部分和空间项部分分别进行泰勒展开得到
          \[
              0 =
              \left(
              T_1 u_t + T_2 u_{tt} + T_3 u_{ttt} + \cdots
              \right)
              +
              \left(
              X_1 u_x + X_2 u_{xx} + X_3 u_{xxx} + \cdots
              \right).
          \]
          将系数对应存储得到下三角矩阵
          \[
              \begin{bmatrix}
                  0                                 \\
                  T_1    & X_1                      \\
                  T_2    &     & X_2                \\
                  T_3    &     &     & X_3          \\
                  \vdots &     &     &     & \ddots \\
              \end{bmatrix},
          \]
          其中 $T_1=1$。

    \item MPDE 方法需要对方程整体关于时间和空间求导，对应系数矩阵的操作为：
          \begin{enumerate}
              \item 关于时间求导表现为系数矩阵整体下移一行；
              \item 关于空间求导表现为系数矩阵整体下移一行再右移一列。
          \end{enumerate}

    \item MPDE 方法的目标是消去时间高阶项和时空混合项，仅保留 $u_t$ 和空间项，对应系数矩阵的操作为逐行从前到后打洞，最终目标是转换为如下形式
          \[
              \begin{bmatrix}
                  0                               \\
                  1 & X_1'                        \\
                    &      & X_2'                 \\
                    &      &      & X_3'          \\
                    &      &      &      & \ddots \\
              \end{bmatrix}.
          \]
          对应
          \[
              0 = u_t + X_1' u_x + X_2' u_{xx} + X_3' u_{xxx} + \cdots.
          \]
\end{enumerate}
具体代码如下（Mathematica 14.0）
\begin{lstlisting}
Init[n_, tlist_, xlist_] :=
  Module[{S, i, j}, S = Table[0, {i, 1, n}, {j, 1, n}];
   For[i = 2, i <= n, i++, S[[i, 1]] = tlist[[i - 1]];
    S[[i, i]] = xlist[[i - 1]]]; S];

Trans[S_, tn_, xn_, n_] :=
 Module[{R, i, j}, R = Table[0, {i, 1, n}, {j, 1, n}];
  For[i = 1, i + tn + xn <= n, i++,
   For[j = 1, j + xn <= n, j++, R[[i + tn + xn, j + xn]] = S[[i, j]]]];
  Return[R]]

Main[S_, n_] :=
 Module[{Cur = S, Tmp},
  For[i = 3, i <= n, i++,
   For[j = 1, j <= i - 1, j++, Tmp = Trans[S, i - j - 1, j - 1, n];
    Cur = Cur - Cur[[i, j]]*Tmp]]; Return[Cur]]

MPDE[n_, tlist_, xlist_] :=
 Module[{S, R}, S = Init[n, tlist, xlist]; R = Main[S, n];
  Return[Simplify[Factor[Diagonal[R]]]]]
\end{lstlisting}

例如考虑对流方程
\[
    u_t + \alpha u_x = 0,\quad (\alpha >0).
\]
以及迎风格式（FTBS格式）
\[
    v_j^{n+1} = v_j^n - \alpha \frac{\Delta t}{\Delta x}(v_j^n - v_j^{n-1}).
\]
基于如下泰勒展开式
\[
    0 =
    \left(
    v_t + v_{tt} \frac{\Delta t}{2!} + v_{ttt} \frac{(\Delta t)^2}{3!} + \cdots
    \right)
    + \alpha
    \left(
    v_x - v_{xx} \frac{\Delta x}{2!} + v_{xxx} \frac{(\Delta x)^2}{3!} + \cdots
    \right).
\]
可以得到两组初始系数
\[
    T_j = \frac{(\Delta t)^{j-1}}{j!}, \quad X_j = \alpha \frac{(-\Delta x)^{j-1}}{j!}, \quad j = 1,2,\dots.
\]
对应的求解代码如下
\begin{lstlisting}
(*u_t + a u_x = 0 FTBS*)
Tlist[n_]:= Table[(dt)^(j-1)/j!,{j,1,n}]
Xlist[n_]:= A * Table[(-dx)^(j-1)/j!,{j,1,n}]
n=5;MPDE[n,Tlist[n],Xlist[n]]
\end{lstlisting}
计算可得系数矩阵的对角线为（记 $r = \alpha \frac{\Delta t}{\Delta x}$）
\[
    0,\quad \alpha,\quad \frac{\alpha \Delta x}2 (r-1),\quad \frac{\alpha (\Delta x)^2}6(r-1)(2r-1),\quad
    \frac{\alpha (\Delta x)^3}{24}(r-1)(6r^2-6r+1).
\]
因此
\begin{align*}
    0 ={} & u_t + \alpha u_x + \frac{\alpha \Delta x}2 (r-1) u_{xx}
    + \frac{\alpha \Delta x^2}6 (r-1)(2r-1) u_{xxx}                   \\
          & + \frac{\alpha (\Delta x)^3}{24}(r-1)(6r^2-6r+1) u_{xxxx}
    + \cdots.
\end{align*}
只需要关注 $u_{xx}$ 和 $u_{xxx}$ 的系数：
\begin{enumerate}
    \item 由于 $u_{xx}$ 的系数为负，差分格式存在耗散性；方程本身无耗散，因此存在数值耗散；
    \item 由于波速为
          \[
              c(k) = \alpha - \frac{\alpha \Delta x^2}6 (r-1)(2r-1) k^2.
          \]
          因此差分格式具有色散性。因为方程的波速为 $c_e(k) = \alpha$，比较可得：
          若 $r \in (\frac12,1)$，差分格式具有数值正色散；若 $r \in (0,\frac12)$，差分格式具有数值负色散。
\end{enumerate}
与前文中通过色散关系推导出的结论一致。


\section{关于 Fourier 方法的推导}

前文中直接给出放大因子并据此检查数值格式的稳定性，显得非常突兀，这里尝试用线性代数以及一点泛函分析的语言进行推导，这种推导方式其实和从 Fourier 理论直接入手的做法殊途同归（参考张强《偏微分方程的有限差分方法》2.4 节）。

\subsection*{周期边界情形}

考虑周期边界下的标量线性常系数单步格式
\[
    \sum_{\ell=-r}^{p} b_\ell v_{j+\ell}^{n+1}
    = \sum_{\ell=-r}^{p} a_\ell v_{j+\ell}^n,
    \quad j=1,\cdots,N.
\]
两端的求和范围可以不同，但是将缺少的系数补零后，总可以写成上面的形式。系数 $a_\ell$ 和 $b_\ell$ 可以依赖于时空网格尺度，但这里假设它们在网格加密过程中保持不变。超过范围的下标在模 $N$ 的意义下理解。
定义循环移位矩阵
\[
    S_N =
    \begin{pmatrix}
        0 & 1 &        &        &   \\
          & 0 & 1      &        &   \\
          &   & \ddots & \ddots &   \\
          &   &        & 0      & 1 \\
        1 &   &        &        & 0
    \end{pmatrix} \in \mathbb{C}^{N\times N}.
\]
定义 Laurent 多项式
\[
    a(z)=\sum_{\ell=-r}^{p}a_\ell z^\ell,
    \qquad
    b(z)=\sum_{\ell=-r}^{p}b_\ell z^\ell.
\]
由于 $S_N$ 可逆，$a(S_N)$ 和 $b(S_N)$ 均有定义。于是差分格式的向量形式为
\[
    b(S_N)V^{n+1}=a(S_N)V^n.
\]
一个合理的单步格式应当保证由 $V^n$ 可以唯一确定 $V^{n+1}$，因此左端算子 $b(S_N)$ 必须可逆。格式于是可以写成
\[
    V^{n+1}=Q_NV^n,
    \qquad
    Q_N=b(S_N)^{-1}a(S_N).
\]
递推可得 $V^n=(Q_N)^nV^0$。

易得 $S_N$ 是酉方阵，自然是正规方阵，可以酉相似对角化
\[
    S_N = W_N D_N W_N^*,
\]
其中 $W_N$ 是酉方阵，$D_N$ 是特征值组成的对角阵，具体为
\[
    D_N = \operatorname{diag}(e^{i\xi_1},\cdots,e^{i\xi_N}), \quad \xi_k=\frac{2\pi k}{N}.
\]
因此 $a(S_N)$ 和 $b(S_N)$ 可以被 $W_N$ 同时酉相似对角化，进而得到 $Q_N$ 也可以酉相似对角化
\[
    Q_N
    = W_N b(D_N)^{-1}a(D_N)W_N^*
    = W_N \operatorname{diag}\left(
    \frac{a(e^{i\xi_1})}{b(e^{i\xi_1})},\dots,
    \frac{a(e^{i\xi_N})}{b(e^{i\xi_N})}
    \right) W_N^*.
\]
因此
\[
    \|(Q_N)^n\| = \max_{k=1,\dots,N} |q(e^{i\xi_k})|^n
\]
其中
\[
    q(z) = \frac{a(z)}{b(z)} = \frac{\sum_{\ell=-r}^{p}a_\ell z^\ell}{\sum_{\ell=-r}^{p}b_\ell z^\ell}.
\]

数值格式的稳定性可以转换为如下要求：存在不依赖 $N$ 的常数 $C$，使得
\[
    \sup_{n\ge 0} \|(Q_N)^n\| \le C.
\]
这里的记号 $\|\,\cdot\,\|$ 是 $L^2$ 范数诱导的算子范数。
对固定的 $N$，要求存在常数 $C_N$ 使得
\[
    \sup_{n\ge 0} \|(Q_N)^n\| \le C_N
\]
等价于要求
\[
    \max_{1\le k \le N} |q(e^{i\xi_k})| \le 1.
\]
但是对于数值格式，我们关注的是网格加密后的极限情况，因此要求有界性对网格加密（$N\to\infty$）一致成立。
由于 $q(z)=a(z)/b(z)$ 在单位圆上有定义且连续，一致稳定性等价于
\[
    \sup_{\xi\in[-\pi,\pi]} |q(e^{i\xi})| \le 1.
\]
这就得到了严格的 von~Neumann 条件，并且
\[
    q(e^{i\xi}) = \frac{a(e^{i\xi})}{b(e^{i\xi})} = \frac{\sum_{\ell=-r}^{p}a_\ell e^{i\ell\xi}}{\sum_{\ell=-r}^{p}b_\ell e^{i\ell\xi}}.
\]
正是之前定义的放大因子。

\subsection*{无界情形}

现在去掉周期边界，考虑整空间网格上的标量线性常系数单步格式
\[
    \sum_{\ell=-r}^{p}b_\ell v_{j+\ell}^{n+1}
    =\sum_{\ell=-r}^{p}a_\ell v_{j+\ell}^n,
    \qquad j\in\mathbb Z.
\]
考虑无限维 Hilbert 空间
\[
    H_h=\ell_h^2(\mathbb Z)
    :=
    \left\{
    V=(V_j)_{j\in\mathbb Z}:
    h\sum_{j\in\mathbb Z}|V_j|^2<\infty
    \right\},
\]
其内积和范数分别为
\[
    \langle U,V\rangle_h
    =
    h\sum_{j\in\mathbb Z}U_j\overline{V_j},
    \qquad
    \|V\|_h
    =
    \left(
    h\sum_{j\in\mathbb Z}|V_j|^2
    \right)^{1/2}.
\]
定义双边移位算子
\[
    S:H_h\to H_h,
    \qquad
    (SV)_j=V_{j+1}.
\]
显然可逆并且逆算子为
\[
    (S^{-1}V)_j=V_{j-1}.
\]
定义算子
\[
    A=a(S),
    \qquad
    B=b(S).
\]
其中 $a(z), b(z)$ 的定义与周期边界情形相同。
单步格式应当保证能够由 $V^n$ 可以唯一确定 $V^{n+1}$，因此左端算子 $B$ 可逆。数值格式可以表示为
\[
    BV^{n+1}=AV^n,
    \qquad
    V^{n+1}=QV^n,
    \qquad
    Q=B^{-1}A,
\]
进而 $V^n=Q^nV^0$。

双边移位算子 $S$ 是酉算子，谱为单位圆
\[
    \sigma(S)
    =
    \left\{
    e^{i \xi}:
    \xi\in[-\pi,\pi]
    \right\}.
\]
由谱映射定理，$Q$ 的谱为
\[
    \sigma(Q)
    =
    q(\sigma(S))
    =
    \left\{
    q(e^{i \xi}):
    \xi\in[-\pi,\pi]
    \right\}.
\]
由于 $S$ 是正规算子，所以 $Q$ 也是正规算子，$Q^n$ 仍然正规，并且
\[
    \|Q^n\| = \max_{\lambda\in\sigma(Q^n)}|\lambda| = \max_{\lambda\in\sigma(Q)}|\lambda|^n = \max_{\xi\in[-\pi,\pi]}|q(e^{i\xi})|^n.
\]
重复之前的稳定性讨论，可以得到无界网格情形下的严格 von Neumann 条件
\[
    \sup_{\xi\in[-\pi,\pi]}
    |q(e^{i\xi})|
    \leq1.
\]
无界情况和周期边界的情况虽然有所区别，但是最终我们得到的是同样的结果。
